% GENERATED FILE. Edit canonical lessons, then run npm run generate:books.
% canonical-source-hash: fnv1a64:f533c2dcfe77d98d
% canonical-lessons: GE-C493-zuneigung, GE-C493-umarmung, GE-C493-zaertlich, GE-C493-geborgenheit, GE-C493-zusammenhalt

\chapter{Zuneigung und Geborgenheit}
\label{ch:ge-zuneigung-und-geborgenheit}
\hlchaptermodality{493}

\begin{chapteropening}
\textbf{By the end of this chapter:} \emph{I can talk about affection and feeling safe.}

Complete the last lesson of chapter 493: I can talk about affection and feeling safe.
\end{chapteropening}

\section[die Zuneigung]{die Zuneigung — affection}

\label{lesson:GE-C493-zuneigung}



\begin{quote}
Before the new one: say the German for to kiss, then the German for kiss.
\end{quote}

\subsection*{You'll want to know: die Zuneigung}
\textbf{die Zuneigung} — \textquotedblleft{}affection\textquotedblright{}.

\begin{quote}
Zwischen den beiden wuchs langsam eine tiefe Zuneigung.
\end{quote}

\begin{quote}
\emph{A deep affection slowly grew between the two.}
\end{quote}

No plural in this sense. Built on \emph{sich zuneigen} — to lean towards.

\begin{cousinweb}
\emph{Zuneigung} is built on \emph{zu} \textquotedblleft{}towards\textquotedblright{} and \emph{neigen} \textquotedblleft{}to lean, to bow\textquotedblright{}, inherited from Middle High German and Old High German \emph{neigen} (EtymDB): leaning towards someone.
\end{cousinweb}

\begin{grammarlens}[title={What you've built}]
One more word for telling, explaining and describing at B1.
\end{grammarlens}

\subsection*{Your turn}
\begin{itemize}
\raggedright
  \item \emph{Say it:} \emph{die Zuneigung}
  \item \emph{Say it:} \emph{die Zuneigung}, once more
  \item \emph{Say it:} a sentence of your own with \emph{die Zuneigung}
  \item \emph{Recall:} say the German for to kiss, then the German for kiss, then say \emph{die Zuneigung} again
\end{itemize}

\begin{tcolorbox}[breakable,colback=teal!4,colframe=teal!35!black,title={Before you move on}]
What does die Zuneigung mean? (\textquotedblleft{}Affection\textquotedblright{}.) Say it once more.
\end{tcolorbox}
\section[die Umarmung]{die Umarmung — hug}

\label{lesson:GE-C493-umarmung}



\begin{quote}
Before the new one: say the German for kiss, then the German for affection.
\end{quote}

\subsection*{You'll want to know: die Umarmung}
\textbf{die Umarmung} — \textquotedblleft{}hug\textquotedblright{}.

\begin{quote}
Zur Begrüßung gab es eine feste Umarmung.
\end{quote}

\begin{quote}
\emph{There was a tight hug by way of hello.}
\end{quote}

Plural: \emph{die Umarmungen}. Built on \emph{umarmen} — to hug.

\begin{cousinweb}
\emph{Umarmung} is built on \emph{umarmen} \textquotedblleft{}to put your arms around\textquotedblright{}, from \emph{um} \textquotedblleft{}around\textquotedblright{} and \emph{Arm}, which is inherited from Old High German \emph{arm}, from Proto-Germanic \emph{armaz} (EtymDB).
\end{cousinweb}

\begin{grammarlens}[title={What you've built}]
One more word for telling, explaining and describing at B1.
\end{grammarlens}

\subsection*{Your turn}
\begin{itemize}
\raggedright
  \item \emph{Say it:} \emph{die Umarmung}
  \item \emph{Say it:} \emph{die Umarmung}, once more
  \item \emph{Say it:} a sentence of your own with \emph{die Umarmung}
  \item \emph{Recall:} say the German for kiss, then the German for affection, then say \emph{die Umarmung} again
\end{itemize}

\begin{tcolorbox}[breakable,colback=teal!4,colframe=teal!35!black,title={Before you move on}]
What does die Umarmung mean? (\textquotedblleft{}Hug\textquotedblright{}.) Say it once more.
\end{tcolorbox}
\section[zärtlich]{zärtlich — tender, affectionate}

\label{lesson:GE-C493-zaertlich}



\begin{quote}
Before the new one: say the German for affection, then the German for hug.
\end{quote}

\subsection*{You'll want to know: zärtlich}
\textbf{zärtlich} — \textquotedblleft{}tender, affectionate\textquotedblright{}.

\begin{quote}
Die Mutter streichelt zärtlich den Kopf des Babys.
\end{quote}

\begin{quote}
\emph{The mother tenderly strokes the baby's head.}
\end{quote}

The noun: \emph{die Zärtlichkeit} — tenderness.

\begin{cousinweb}
\emph{zärtlich} is built on \emph{zart} \textquotedblleft{}delicate, tender\textquotedblright{}, inherited from Old High German \emph{zart}, from Proto-Germanic \emph{tardaz} (EtymDB), with \emph{-lich}.
\end{cousinweb}

\begin{grammarlens}[title={What you've built}]
One more word for telling, explaining and describing at B1.
\end{grammarlens}

\subsection*{Your turn}
\begin{itemize}
\raggedright
  \item \emph{Say it:} \emph{zärtlich}
  \item \emph{Say it:} \emph{zärtlich}, once more
  \item \emph{Say it:} a sentence of your own with \emph{zärtlich}
  \item \emph{Recall:} say the German for affection, then the German for hug, then say \emph{zärtlich} again
\end{itemize}

\begin{tcolorbox}[breakable,colback=teal!4,colframe=teal!35!black,title={Before you move on}]
What does zärtlich mean? (\textquotedblleft{}Tender, affectionate\textquotedblright{}.) Say it once more.
\end{tcolorbox}
\section[die Geborgenheit]{die Geborgenheit — feeling of safety, security}

\label{lesson:GE-C493-geborgenheit}



\begin{quote}
Before the new one: say the German for hug, then the German for tender.
\end{quote}

\subsection*{You'll want to know: die Geborgenheit}
\textbf{die Geborgenheit} — \textquotedblleft{}feeling of safety, security\textquotedblright{}.

\begin{quote}
Bei meinen Großeltern spürte ich immer Geborgenheit.
\end{quote}

\begin{quote}
\emph{At my grandparents' I always felt safe and secure.}
\end{quote}

No plural. Built on \emph{geborgen} — safe, sheltered.

\begin{cousinweb}
\emph{Geborgenheit} is built on \emph{geborgen} \textquotedblleft{}sheltered\textquotedblright{}, the participle of \emph{bergen} \textquotedblleft{}to shelter, to save\textquotedblright{}, inherited from Old High German \emph{bergan}, from Proto-Germanic \emph{berganą} (EtymDB).
\end{cousinweb}

\begin{grammarlens}[title={What you've built}]
One more word for telling, explaining and describing at B1.
\end{grammarlens}

\subsection*{Your turn}
\begin{itemize}
\raggedright
  \item \emph{Say it:} \emph{die Geborgenheit}
  \item \emph{Say it:} \emph{die Geborgenheit}, once more
  \item \emph{Say it:} a sentence of your own with \emph{die Geborgenheit}
  \item \emph{Recall:} say the German for hug, then the German for tender, then say \emph{die Geborgenheit} again
\end{itemize}

\begin{tcolorbox}[breakable,colback=teal!4,colframe=teal!35!black,title={Before you move on}]
What does die Geborgenheit mean? (\textquotedblleft{}Feeling of safety, security\textquotedblright{}.) Say it once more.
\end{tcolorbox}
\section[der Zusammenhalt]{der Zusammenhalt — solidarity, sticking together}

\label{lesson:GE-C493-zusammenhalt}



\begin{quote}
Before the new one: say the German for tender, then the German for feeling of safety.
\end{quote}

\subsection*{You'll want to know: der Zusammenhalt}
\textbf{der Zusammenhalt} — \textquotedblleft{}solidarity, sticking together\textquotedblright{}.

\begin{quote}
In unserer Familie ist der Zusammenhalt sehr stark.
\end{quote}

\begin{quote}
\emph{In our family we really stick together.}
\end{quote}

No plural. Built on \emph{zusammenhalten} — to stick together.

\begin{cousinweb}
\emph{Zusammenhalt} is built on \emph{zusammen} \textquotedblleft{}together\textquotedblright{} and \emph{Halt} \textquotedblleft{}hold\textquotedblright{}, from \emph{halten} \textquotedblleft{}to hold\textquotedblright{} (EtymDB), a Germanic verb from Proto-Germanic \emph{haldaną}.
\end{cousinweb}

\begin{grammarlens}[title={What you've built}]
One more word for telling, explaining and describing at B1.
\end{grammarlens}

\subsection*{Your turn}
\begin{itemize}
\raggedright
  \item \emph{Say it:} \emph{der Zusammenhalt}
  \item \emph{Say it:} \emph{der Zusammenhalt}, once more
  \item \emph{Say it:} a sentence of your own with \emph{der Zusammenhalt}
  \item \emph{Recall:} say the German for tender, then the German for feeling of safety, then say \emph{der Zusammenhalt} again
\end{itemize}

\begin{tcolorbox}[breakable,colback=teal!4,colframe=teal!35!black,title={Before you move on}]
What does der Zusammenhalt mean? (\textquotedblleft{}Solidarity, sticking together\textquotedblright{}.) Say it once more.
\end{tcolorbox}
